\chapter{测量}\label{ch:g3:measure}

有多长？有多重？能装多少？几点了？测量
能把世界变成数字 --- 前提是大家使用相同的单位。
本章介绍长度、质量和容积的单位，
并教你如何读钟表。

\section{长度}

\begin{definition}[长度单位]\label{def:g3:measure:length}
长度的基本单位是\emph{米}（m）。较小的长度用
厘米（cm）和毫米（mm）；较长的距离用
千米（km）。
\[
1 \text{ 米} = 100 \text{ 厘米}, \qquad
1 \text{ 厘米} = 10 \text{ 毫米}, \qquad
1 \text{ 千米} = 1\,000 \text{ 米}.
\]
\end{definition}

\begin{omfigure}
\begin{tikzpicture}[scale=1.3]
  % ruler body extends past the marks so no numeral is clipped
  \draw[thick] (-0.35,0) rectangle (8.35,0.9);
  \foreach \i in {0,...,8} {
    \draw[thick] (\i,0.9) -- (\i,0.55);
    \node[font=\small] at (\i,0.32) {\i};
  }
  \foreach \i in {0,...,80} \draw (\i*0.1,0.9) -- (\i*0.1,0.72);
  \draw[omThm, very thick] (0,1.25) -- (3.6,1.25);
  \node[omThm, above, font=\small] at (1.8,1.3)
    {$3.6$ 厘米 $= 3$ 厘米和 $6$ 毫米};
\end{tikzpicture}

{\small 读尺子：大刻度表示厘米，小刻度表示
毫米。物体应从 $0$ 刻度开始 --- 而不是从
尺子的边缘开始！ --- 终点在第 $6$ 个小刻度，位于 $3$ 之后。}
\end{omfigure}

\begin{example}[选择合适的单位]\label{ex:g3:measure:choose}
一只蚂蚁：约 $5$ 毫米。一支铅笔：约 $18$ 厘米。一扇门：约 $2$ 米。
走到下一个村庄：约 $3$ 千米。说“铅笔的长度是
$18$”毫无意义 --- 单位是答案的一部分。
\end{example}

\begin{method}[单位换算]\label{met:g3:measure:convert}
换算时，乘以或除以 $10$、$100$ 或 $1\,000$：
\begin{enumerate}
  \item 由大单位换成小单位，\emph{乘}：
        $3$ 米 $= 3 \times 100 = 300$ 厘米；
  \item 由小单位换成大单位，\emph{除}：
        $250$ 厘米 $= 250 \div 100 = 2$ 米 和 $50$ 厘米；
  \item 要把长度相加，须先将它们换算成\emph{相同}的单位：
        $1$ 米 $+ 35$ 厘米 $= 100 + 35 = 135$ 厘米。
\end{enumerate}
\end{method}

\section{质量和容积}

\begin{definition}[质量和容积的单位]\label{def:g3:measure:mass}
质量用\emph{克}（g）和\emph{千克}（kg）计量，
容积（即容器所能容纳的量）用\emph{升}（L）和
\emph{厘升}（cL）计量：
\[
1 \text{ 千克} = 1\,000 \text{ 克},
\qquad
1 \text{ 升} = 100 \text{ 厘升}.
\]
\end{definition}

\begin{example}\label{ex:g3:measure:mass}
一封信：约 $20$ 克。一袋面粉：$1$ 千克。一个浴缸：约
$150$ 升。一只玻璃杯：约 $20$ 厘升。要使秤的一端放着一个
$1$ 千克的袋子，另一端放 $650$ 克苹果，需在
苹果一侧加上 $1\,000 - 650 = 350$ 克。
\end{example}

\section{认读钟表}

\begin{method}[读时间]\label{met:g3:measure:clock}
对于带指针的钟表：
\begin{enumerate}
  \item 短的\emph{时针}指示小时：读出它刚刚
        经过的数字；
  \item 长的\emph{分针}指示分钟：每个大格为 $5$
        分钟，所以分针指向 $3$ 时表示 $15$ 分钟；
  \item 过了半点之后，还可以读作“差……到下一个整点”：
        $10{:}50$ 读作“差十分十一点”。
\end{enumerate}
\end{method}

\begin{omfigure}
\begin{tikzpicture}[scale=1.0]
  % clock 1: 3:00
  \draw[thick] (0,0) circle (1.3);
  \foreach \h in {1,...,12} {
    \node[font=\footnotesize] at ({90-\h*30}:1.05) {\h};
    \draw ({90-\h*30}:1.22) -- ({90-\h*30}:1.3);
  }
  \draw[very thick, omDef] (0,0) -- (0:0.65);
  \draw[thick, omThm] (0,0) -- (90:1.05);
  \node[below, font=\small] at (0,-1.55) {3:00};
  % clock 2: 7:30
  \begin{scope}[xshift=4.4cm]
  \draw[thick] (0,0) circle (1.3);
  \foreach \h in {1,...,12} {
    \node[font=\footnotesize] at ({90-\h*30}:1.05) {\h};
    \draw ({90-\h*30}:1.22) -- ({90-\h*30}:1.3);
  }
  \draw[very thick, omDef] (0,0) -- ({90-7.5*30}:0.65);
  \draw[thick, omThm] (0,0) -- (270:1.05);
  \node[below, font=\small] at (0,-1.55) {7:30};
  \end{scope}
  % clock 3: 10:50
  \begin{scope}[xshift=8.8cm]
  \draw[thick] (0,0) circle (1.3);
  \foreach \h in {1,...,12} {
    \node[font=\footnotesize] at ({90-\h*30}:1.05) {\h};
    \draw ({90-\h*30}:1.22) -- ({90-\h*30}:1.3);
  }
  \draw[very thick, omDef] (0,0) -- ({90-10.83*30}:0.65);
  \draw[thick, omThm] (0,0) -- ({90-10*30}:1.05);
  \node[below, font=\small] at (0,-1.55) {10:50};
  \end{scope}
\end{tikzpicture}

{\small 三只钟表：短的蓝色时针表示时，长的红色分针表示分。7:30 时，
时针位于 7 和 8 \emph{之间}；10:50 时，
时针几乎指向 11。}
\end{omfigure}

\begin{example}[时长]\label{ex:g3:measure:duration}
影片 $14{:}40$ 开始，$16{:}10$ 结束。它放映了多长时间？
分两步数一数：
\begin{enumerate}
  \item 从 $14{:}40$ 到 $15{:}00$：$20$ 分钟；
  \item 从 $15{:}00$ 到 $16{:}10$：$1$ 小时 $10$ 分钟。
\end{enumerate}
合计：$1$ 小时 $30$ 分钟。注意：一小时有 $60$ 分钟，所以
$16{:}10$ 减 $14{:}40$ 不能像 $1610 - 1440$ 那样计算！
\end{example}

\section{练习}

\begin{exercise}[$\star$]\label{exo:g3:measure:1}
下面各量用哪个单位（毫米、厘米、米或千米）比较合适：钉子的宽度、
同学的身高、学校操场的长度、
两个城镇之间的距离？
\end{exercise}

\begin{exercise}[$\star$]\label{exo:g3:measure:2}
换算：$5$ 米 = ？厘米；\quad $70$ 毫米 = ？厘米；\quad $4$ 千米 = ？米；
\quad $600$ 厘米 = ？米。
\end{exercise}

\begin{exercise}[$\star$]\label{exo:g3:measure:3}
用直尺（精确到毫米）量出同学所画三角形各边的长度，
并计算它的
周长。
\end{exercise}

\begin{exercise}[$\star$]\label{exo:g3:measure:4}
先换算，再计算：$2$ 米 $+ 45$ 厘米（用厘米表示）；\quad
$1$ 千米 $+ 250$ 米（用米表示）；\quad $3$ 厘米 $+ 8$ 毫米（用毫米表示）。
\end{exercise}

\begin{exercise}[$\star$]\label{exo:g3:measure:5}
换算：$3$ 千克 = ？克；\quad $2\,500$ 克 = ？千克 ？克；\quad
$4$ 升 = ？厘升；\quad $350$ 厘升 = ？升 ？厘升。
\end{exercise}

\begin{exercise}[$\star$]\label{exo:g3:measure:6}
做一份食谱需要 $750$ 克面粉。莫娜有一袋 $1$ 千克的面粉，
够用吗？还剩多少？
\end{exercise}

\begin{exercise}[$\star$]\label{exo:g3:measure:7}
读出下列时刻：时针在 $2$ 和 $3$ 之间，分针指向 $6$；
时针指向 $9$，分针指向 $12$；
时针几乎指向 $5$，分针指向 $10$。
\end{exercise}

\begin{exercise}[$\star$]\label{exo:g3:measure:8}
学校 $8{:}30$ 上课，上午在 $11{:}45$ 结束。上午持续了
多长时间？（像
\cref{ex:g3:measure:duration}那样分步数一数。）
\end{exercise}

\begin{exercise}[$\star$]\label{exo:g3:measure:9}
火车 $9{:}50$ 出发，行程持续 $45$ 分钟，
它几点到达？（先数到 $10{:}00$。）
\end{exercise}

\begin{exercise}[$\star\star$]\label{exo:g3:measure:10}
一个瓶子能装 $1$ 升水。保罗要把水倒进每杯 $20$ 厘升的杯子里，
他最多能倒满几杯？（先换算。）
\end{exercise}

\begin{exercise}[$\star\star$]\label{exo:g3:measure:11}
莉娜每天走 $600$ 米去学校，再走回家。
一周上学 $5$ 天，她一共走多少米？
请先用米作答，再换算成千米。
\end{exercise}
